% GENERATED FILE. Edit canonical lessons, then run npm run generate:books.
% canonical-source-hash: fnv1a64:a09475a14fa7603b
% canonical-lessons: MR-C316-lagna, MR-C316-sakharpuda, MR-C316-mendi, MR-C316-varat, MR-C316-mandap

\chapter{A Wedding}
\label{ch:mr-a-wedding}
\hlchaptermodality{316}

\begin{chapteropening}
\textbf{By the end of this chapter:} \emph{I can talk about a wedding: the engagement, the henna, the procession.}

Complete the last lesson of chapter 316: I can talk about a wedding: the engagement, the henna, the procession.
\end{chapteropening}

\section[\texorpdfstring{\mr{लग्न} (lagna)}{लग्न (lagna)}]{\mr{लग्न} (lagna) — a wedding; a marriage}

\label{lesson:MR-C316-lagna}



\begin{quote}
Before the new one: say the Marathi for to register, then the Marathi for to wander.
\end{quote}

\subsection*{You'll want to know: \mr{लग्न}}
\textbf{\mr{लग्न}} — \emph{lagna} — \textquotedblleft{}a wedding; a marriage\textquotedblright{}.

\begin{quote}
\mr{माझ्या} \mr{मावशीचं} \mr{लग्न} \mr{पुढच्या} \mr{महिन्यात} \mr{आहे}.
\end{quote}

\begin{quote}
\emph{mājhyā māvśīcaṁ lagna puḍhacyā mahinyāt āhe.}
\end{quote}

\begin{quote}
My aunt's wedding is next month.
\end{quote}

Neuter.

\mr{लग्न} \mr{होणे}, to get married; X-\mr{शी} \mr{लग्न} \mr{करणे}, to marry X.

\begin{cousinweb}
\textbf{\mr{लग्न}} (\emph{lagna}) is a tatsama, taken from Sanskrit \emph{lagna}, 'attached, joined, fixed' (Monier-Williams): a joining of two people and two families.
\end{cousinweb}

\begin{grammarlens}[title={What you've built}]
One more word for telling what happened.
\end{grammarlens}

\subsection*{Your turn}
\begin{itemize}
\raggedright
  \item \emph{Say it:} \emph{lagna}
  \item \emph{Say it:} \emph{lagna}, once more
  \item \emph{Say it:} \emph{mājhyā māvśīcaṁ lagna puḍhacyā mahinyāt āhe.} — My aunt's wedding is next month.
  \item \emph{Recall:} say the Marathi for to register, then the Marathi for to wander, then say \emph{lagna} again
\end{itemize}

\begin{tcolorbox}[breakable,colback=teal!4,colframe=teal!35!black,title={Before you move on}]
What does \mr{लग्न} mean? (\textquotedblleft{}A wedding; a marriage\textquotedblright{}.) Say it once more.
\end{tcolorbox}
\section[\texorpdfstring{\mr{साखरपुडा} (sākharpuḍā)}{साखरपुडा (sākharpuḍā)}]{\mr{साखरपुडा} (sākharpuḍā) — an engagement}

\label{lesson:MR-C316-sakharpuda}



\begin{quote}
Before the new one: say the Marathi for to wander, then the Marathi for a wedding.
\end{quote}

\subsection*{You'll want to know: \mr{साखरपुडा}}
\textbf{\mr{साखरपुडा}} — \emph{sākharpuḍā} — \textquotedblleft{}an engagement\textquotedblright{}.

\begin{quote}
\mr{त्यांचा} \mr{साखरपुडा} \mr{काल} \mr{घरीच} \mr{झाला}.
\end{quote}

\begin{quote}
\emph{tyāñcā sākharpuḍā kāl gharīc jhālā.}
\end{quote}

\begin{quote}
Their engagement took place at home yesterday.
\end{quote}

Masculine.

A packet of sugar is exchanged: \mr{साखर} + \mr{पुडा}.

\begin{cousinweb}
\textbf{\mr{साखरपुडा}} (\emph{sākharpuḍā}) is built from \emph{sākhar}, 'sugar', inherited from Sanskrit \emph{śarkarā} (CDIAL 12338), and \emph{puḍā}, 'a packet', which CDIAL lists under Sanskrit \emph{puṭa}, 'a fold, a small container' (CDIAL 8253).
\end{cousinweb}

\begin{grammarlens}[title={What you've built}]
One more word for telling what happened.
\end{grammarlens}

\subsection*{Your turn}
\begin{itemize}
\raggedright
  \item \emph{Say it:} \emph{sākharpuḍā}
  \item \emph{Say it:} \emph{sākharpuḍā}, once more
  \item \emph{Say it:} \emph{tyāñcā sākharpuḍā kāl gharīc jhālā.} — Their engagement took place at home yesterday.
  \item \emph{Recall:} say the Marathi for to wander, then the Marathi for a wedding, then say \emph{sākharpuḍā} again
\end{itemize}

\begin{tcolorbox}[breakable,colback=teal!4,colframe=teal!35!black,title={Before you move on}]
What does \mr{साखरपुडा} mean? (\textquotedblleft{}An engagement\textquotedblright{}.) Say it once more.
\end{tcolorbox}
\section[\texorpdfstring{\mr{मेंदी} (mendī)}{मेंदी (mendī)}]{\mr{मेंदी} (mendī) — henna}

\label{lesson:MR-C316-mendi}



\begin{quote}
Before the new one: say the Marathi for a wedding, then the Marathi for an engagement.
\end{quote}

\subsection*{You'll want to know: \mr{मेंदी}}
\textbf{\mr{मेंदी}} — \emph{mendī} — \textquotedblleft{}henna\textquotedblright{}.

\begin{quote}
\mr{लग्नाआधी} \mr{नवरीच्या} \mr{हातांवर} \mr{मेंदी} \mr{काढतात}.
\end{quote}

\begin{quote}
\emph{lagnāādhī navrīcyā hātāṁvar mendī kāḍhtāt.}
\end{quote}

\begin{quote}
Before the wedding, henna is put on the bride's hands.
\end{quote}

Feminine.

\mr{मेंदी} \mr{काढणे}, to put on henna in patterns.

\begin{cousinweb}
\textbf{\mr{मेंदी}} (\emph{mendī}) is a tadbhava: CDIAL lists Marathi \emph{mẽdhī}, 'the henna plant', under Sanskrit \emph{mendhī} or \emph{mendhikā}, 'Lawsonia alba' (CDIAL 10328).
\end{cousinweb}

\begin{grammarlens}[title={What you've built}]
One more word for telling what happened.
\end{grammarlens}

\subsection*{Your turn}
\begin{itemize}
\raggedright
  \item \emph{Say it:} \emph{mendī}
  \item \emph{Say it:} \emph{mendī}, once more
  \item \emph{Say it:} \emph{lagnāādhī navrīcyā hātāṁvar mendī kāḍhtāt.} — Before the wedding, henna is put on the bride's hands.
  \item \emph{Recall:} say the Marathi for a wedding, then the Marathi for an engagement, then say \emph{mendī} again
\end{itemize}

\begin{tcolorbox}[breakable,colback=teal!4,colframe=teal!35!black,title={Before you move on}]
What does \mr{मेंदी} mean? (\textquotedblleft{}Henna\textquotedblright{}.) Say it once more.
\end{tcolorbox}
\section[\texorpdfstring{\mr{वरात} (varāt)}{वरात (varāt)}]{\mr{वरात} (varāt) — a wedding procession}

\label{lesson:MR-C316-varat}



\begin{quote}
Before the new one: say the Marathi for an engagement, then the Marathi for henna.
\end{quote}

\subsection*{You'll want to know: \mr{वरात}}
\textbf{\mr{वरात}} — \emph{varāt} — \textquotedblleft{}a wedding procession\textquotedblright{}.

\begin{quote}
\mr{वरात} \mr{वाजत} \mr{गाजत} \mr{मंडपात} \mr{पोहोचली}.
\end{quote}

\begin{quote}
\emph{varāt vājat gājat maṇḍpāt pohoclī.}
\end{quote}

\begin{quote}
The wedding procession reached the marquee with music and fanfare.
\end{quote}

Feminine.

The groom and his party come in it.

\begin{cousinweb}
\textbf{\mr{वरात}} (\emph{varāt}) is a tadbhava: CDIAL lists Marathi \emph{varāt}, 'the groom's procession', under Sanskrit \emph{varayātrā}, 'the procession of a suitor or bridegroom' (CDIAL 11322).
\end{cousinweb}

\begin{grammarlens}[title={What you've built}]
One more word for telling what happened.
\end{grammarlens}

\subsection*{Your turn}
\begin{itemize}
\raggedright
  \item \emph{Say it:} \emph{varāt}
  \item \emph{Say it:} \emph{varāt}, once more
  \item \emph{Say it:} \emph{varāt vājat gājat maṇḍpāt pohoclī.} — The wedding procession reached the marquee with music and fanfare.
  \item \emph{Recall:} say the Marathi for an engagement, then the Marathi for henna, then say \emph{varāt} again
\end{itemize}

\begin{tcolorbox}[breakable,colback=teal!4,colframe=teal!35!black,title={Before you move on}]
What does \mr{वरात} mean? (\textquotedblleft{}A wedding procession\textquotedblright{}.) Say it once more.
\end{tcolorbox}
\section[\texorpdfstring{\mr{मंडप} (maṇḍap)}{मंडप (maṇḍap)}]{\mr{मंडप} (maṇḍap) — a marquee; a pavilion}

\label{lesson:MR-C316-mandap}



\begin{quote}
Before the new one: say the Marathi for henna, then the Marathi for a wedding procession.
\end{quote}

\subsection*{You'll want to know: \mr{मंडप}}
\textbf{\mr{मंडप}} — \emph{maṇḍap} — \textquotedblleft{}a marquee; a pavilion\textquotedblright{}.

\begin{quote}
\mr{अंगणात} \mr{लग्नासाठी} \mr{मोठा} \mr{मंडप} \mr{उभा} \mr{केला}.
\end{quote}

\begin{quote}
\emph{aṅgṇāt lagnāsāṭhī moṭhā maṇḍap ubhā kelā.}
\end{quote}

\begin{quote}
A big marquee was put up in the yard for the wedding.
\end{quote}

Masculine.

\mr{मंडप} \mr{उभा} \mr{करणे} / \mr{घालणे}, to put up a marquee.

\begin{cousinweb}
\textbf{\mr{मंडप}} (\emph{maṇḍap}) is a tatsama, taken from Sanskrit \emph{maṇḍapa}, 'an open hall, a temporary pavilion' (Monier-Williams).
\end{cousinweb}

\begin{grammarlens}[title={What you've built}]
One more word for telling what happened.
\end{grammarlens}

\subsection*{Your turn}
\begin{itemize}
\raggedright
  \item \emph{Say it:} \emph{maṇḍap}
  \item \emph{Say it:} \emph{maṇḍap}, once more
  \item \emph{Say it:} \emph{aṅgṇāt lagnāsāṭhī moṭhā maṇḍap ubhā kelā.} — A big marquee was put up in the yard for the wedding.
  \item \emph{Recall:} say the Marathi for henna, then the Marathi for a wedding procession, then say \emph{maṇḍap} again
\end{itemize}

\begin{tcolorbox}[breakable,colback=teal!4,colframe=teal!35!black,title={Before you move on}]
What does \mr{मंडप} mean? (\textquotedblleft{}A marquee; a pavilion\textquotedblright{}.) Say it once more.
\end{tcolorbox}
